\origpage{163}{GRAPH MINORS, XVII}
\noindent which every edge has two or three ends. To carry out the encoding procedure, however, we need to arrange that the near-embeddings of the pieces have nice connectivity properties. The object of this paper is to show that any graph which can be near-embedded in some surface can also be ``nicely'' near-embedded in some surface, homeomorphic to or simpler than the first, after removing a bounded number of vertices. (Readers are warned that, unless they understand Theorem (3.1) of [2], further reading of the present paper is a waste of time.) Most of the paper is concerned with, given a near-embedding, finding a better one, and analyzing its structure when there is no better one. In Section 13 we connect this with excluding a minor.

The main result of this paper is therefore rather humdrum; it will not cause the reader any excitement, and its proof is unfortunately quite lengthy. To stimulate the reader's interest, let us mention that \textit{finding} the result was not humdrum at all. We really needed this to be true, for all the applications of this series of papers, and for a long time the proof eluded us. (The crucial idea was that of ``warp,'' defined in Section 3. This may seem unnatural, but nothing simpler works as far as we can see.)

Since for some applications it is needed, we work with near-embeddings of hypergraphs rather than of graphs. Near-embeddings of hypergraphs arise when a minor is excluded from the ``1-skeleton'' of the hypergraph, a situation which occurs in applications. This is discussed in Sections 12 and 14. So a second objective of the paper is to convert near-embeddings of 1-skeletons of hypergraphs to near-embeddings of the hypergraphs themselves. This is straightforward, however, and is nowhere near as delicate or difficult as the first objective.

A surface is a compact 2-manifold with (possibly empty) boundary. We denote the boundary of a surface $\Sigma$ by $bd(\Sigma)$. Each component of $bd(\Sigma)$ is homeomorphic to a circle, and we call these components the \textit{cuffs} of $\Sigma$. We denote by $\Sigma(a, b, c)$ a surface obtained from a sphere by adding $a$ handles and $b$ crosscaps and deleting the interiors of $c$ mutually disjoint closed discs. Thus $\Sigma(0, 0, 1)$ is a closed disc (which we shall usually just call a \textit{disc}; when we mean an open disc we shall say so), $\Sigma(0, 0, 2)$ is a cylinder, $\Sigma(1, 0, 0)$ is a torus, etc. Surfaces in this paper will usually be connected. It is known that

\stmt
(1.1) (i) \textit{For every non-null connected surface $\Sigma$ there are integers $a$, $b$, $c$ such that $\Sigma \simeq \Sigma(a, b, c)$;}

\hspace*{12pt}(ii) \textit{$\Sigma(a, b, c) \simeq \Sigma(a', b', c')$ if and only if $c = c'$, $\delm{a + 2b = a' + 2b'}\allowbreak\fix{2a + b = 2a' + b'}$\fixnote{fix:1.1}{The coefficients of $a$ (handles) and $b$ (crosscaps) are swapped. For example, $\Sigma(1, 1, 0) \simeq \Sigma(0, 3, 0)$ by Dyck's theorem, and $2a + b = 3 = 2a' + b'$, but $a + 2b = 3 \neq 6 = a' + 2b'$. Condition (ii) of ``simpler than'' on page 164 uses $2a + b$. With the correction, (1.1) holds.}, and both or neither of $b$, $b'$ are zero.}

[ $\simeq$ denotes homeomorphism. ]
